\documentclass[10pt,a4paper]{amsart}
\usepackage{amsmath,amssymb,amsthm,graphicx}
\usepackage{verbatim}
\newenvironment{reference}{\comment}{\endcomment}
%\newenvironment{reference}{}{}
\newenvironment{slogan}{\comment}{\endcomment}
\newenvironment{history}{\comment}{\endcomment}

% For commutative diagrams we use Xy-pic
\usepackage[all]{xy}

% We use 2cell for 2-commutative diagrams.
\xyoption{2cell}
% Theorem environments.
%
\theoremstyle{plain}
\newtheorem{theorem}[subsection]{Theorem}
\newtheorem{proposition}[subsection]{Proposition}
\newtheorem{lemma}[subsection]{Lemma}

\theoremstyle{definition}
\newtheorem{definition}[subsection]{Definition}
\newtheorem{example}[subsection]{Example}
\newtheorem{exercise}[subsection]{Exercise}
\newtheorem{situation}[subsection]{Situation}

\theoremstyle{remark}
\newtheorem{remark}[subsection]{Remark}
\newtheorem{remarks}[subsection]{Remarks}

\numberwithin{equation}{subsection}

% Macros
%
\def\lim{\mathop{\mathrm{lim}}\nolimits}
\def\colim{\mathop{\mathrm{colim}}\nolimits}
\def\Spec{\mathop{\mathrm{Spec}}}
\def\Hom{\mathop{\mathrm{Hom}}\nolimits}
\def\Ext{\mathop{\mathrm{Ext}}\nolimits}
\def\SheafHom{\mathop{\mathcal{H}\!\mathit{om}}\nolimits}
\def\SheafExt{\mathop{\mathcal{E}\!\mathit{xt}}\nolimits}
\def\Sch{\mathit{Sch}}
\def\Mor{\mathop{\mathrm{Mor}}\nolimits}
\def\Ob{\mathop{\mathrm{Ob}}\nolimits}
\def\Sh{\mathop{\mathit{Sh}}\nolimits}
\def\NL{\mathop{N\!L}\nolimits}
\def\CH{\mathop{\mathrm{CH}}\nolimits}
\def\proetale{{pro\text{-}\acute{e}tale}}
\def\etale{{\acute{e}tale}}
\def\QCoh{\mathit{QCoh}}
\def\Ker{\mathop{\mathrm{Ker}}}
\def\Im{\mathop{\mathrm{Im}}}
\def\Coker{\mathop{\mathrm{Coker}}}
\def\Coim{\mathop{\mathrm{Coim}}}

% Boxtimes
%
\DeclareMathSymbol{\boxtimes}{\mathbin}{AMSa}{"02}

%
% Macros for moduli stacks/spaces
%
\def\QCohstack{\mathcal{QC}\!\mathit{oh}}
\def\Cohstack{\mathcal{C}\!\mathit{oh}}
\def\Spacesstack{\mathcal{S}\!\mathit{paces}}
\def\Quotfunctor{\mathrm{Quot}}
\def\Hilbfunctor{\mathrm{Hilb}}
\def\Curvesstack{\mathcal{C}\!\mathit{urves}}
\def\Polarizedstack{\mathcal{P}\!\mathit{olarized}}
\def\Complexesstack{\mathcal{C}\!\mathit{omplexes}}
% \Pic is the operator that assigns to X its picard group, usage \Pic(X)
% \Picardstack_{X/B} denotes the Picard stack of X over B
% \Picardfunctor_{X/B} denotes the Picard functor of X over B
\def\Pic{\mathop{\mathrm{Pic}}\nolimits}
\def\Picardstack{\mathcal{P}\!\mathit{ic}}
\def\Picardfunctor{\mathrm{Pic}}
\def\Deformationcategory{\mathcal{D}\!\mathit{ef}}
\usepackage{cleveref}
\crefname{originaldefinition}{original definition}{original definitions}
\Crefname{originaldefinition}{Original definition}{Original definitions}
\makeatletter
\newlabel{categories-definition-directed}{{002V}{source}{}{sourcepointer-1}{}}
\newlabel{categories-definition-directed@cref}{{[originaldefinition][0][]002V}{[1][0][]source}{}{}{}}
\makeatother
\crefname{originalsection}{original section}{original sections}
\Crefname{originalsection}{Original section}{Original sections}
\makeatletter
\newlabel{categories-section-limits}{{002D}{source}{}{sourcepointer-2}{}}
\newlabel{categories-section-limits@cref}{{[originalsection][0][]002D}{[1][0][]source}{}{}{}}
\makeatother
\crefname{originalreference}{original reference}{original references}
\Crefname{originalreference}{Original reference}{Original references}
\makeatletter
\newlabel{lemma-extend}{{01PF}{source}{}{sourcepointer-3}{}}
\newlabel{lemma-extend@cref}{{[originalreference][0][]01PF}{[1][0][]source}{}{}{}}
\makeatother
\crefname{originalreference}{original reference}{original references}
\Crefname{originalreference}{Original reference}{Original references}
\makeatletter
\newlabel{lemma-extend-finite-presentation}{{01PI}{source}{}{sourcepointer-4}{}}
\newlabel{lemma-extend-finite-presentation@cref}{{[originalreference][0][]01PI}{[1][0][]source}{}{}{}}
\makeatother
\crefname{originallemma}{original lemma}{original lemmas}
\Crefname{originallemma}{Original lemma}{Original lemmas}
\makeatletter
\newlabel{modules-lemma-set-isomorphism-classes-finite-type-modules}{{01BC}{source}{}{sourcepointer-5}{}}
\newlabel{modules-lemma-set-isomorphism-classes-finite-type-modules@cref}{{[originallemma][0][]01BC}{[1][0][]source}{}{}{}}
\makeatother
\crefname{originalsection}{original section}{original sections}
\Crefname{originalsection}{Original section}{Original sections}
\makeatletter
\newlabel{sheaves-section-limits-sheaves}{{009E}{source}{}{sourcepointer-6}{}}
\newlabel{sheaves-section-limits-sheaves@cref}{{[originalsection][0][]009E}{[1][0][]source}{}{}{}}
\makeatother
\title{A quasi-coherent module on a quasi-compact and quasi-separated scheme is a filtered colimit of finitely presented modules}
\author{The Stacks Project Authors}
\date{stacks-project:properties.tex:local-snapshot-20261012; GNU FDL 1.2 or later}
\begin{document}
\maketitle
\noindent\textit{Editorial note.} All delivered mathematical text is the author's own, copied from the pinned Stacks Project chapter file properties.tex; no statement or proof step was rewritten, reordered or removed.\par\medskip
\noindent\textit{Editorial note.} The counted proof is the author's complete argument: the enumeration of the finitely presented modules mapping to the given one, the verification of filteredness through direct sums and cokernels of differences, and the stalkwise comparison of the colimit with the module using the two extension lemmas.\par\medskip
\noindent\textit{Editorial note.} The label lemma-directed-colimit-diagram-finite-presentation names the counted statement, and the proof environment immediately following it in the source is delivered verbatim as its proof.\par\medskip
\noindent\textit{Editorial note.} The counted statement is carried by the author's own lemma environment in the chapter source.\par\medskip
\noindent\textit{Editorial note.} Pre-open gate: python operations/agent-authoring/tools/precard\_gate.py --source properties.tex --label lemma-directed-colimit-diagram-finite-presentation --start 3256 --end 3407 --local-id 28041 returned RESULT=OK (L28041 free, occupancy FREE, public interval clear). Candidate taken from sources/stacks-properties/candidates-halfA.json.\par\medskip
\noindent\textit{Editorial note.} The statement opens with a native slogan environment stating the content informally; that environment is part of the delivered statement and is covered by the preamble span 8-13 of the chapter, which this card includes.\par\medskip
\noindent\textit{Editorial note.} The statement and proof are copied from the pinned chapter file; the cited inputs are kept as external pointers with their canonical Stacks tags.\par\medskip
\noindent\textit{Allowed standard prerequisites (editorial).}\begin{itemize}
\item Boundedness of the isomorphism classes of finitely presented modules. Modules, the lemma that the finitely presented modules on a scheme form a set of isomorphism classes bounded by a cardinal. Supplies the existence of the index set I enumerating the finitely presented modules mapping to the given module.
\item Extension of a finite type subsheaf from a quasi-compact open. Same chapter, the lemma producing a finite type quasi-coherent subsheaf restricting to a given one, and the corresponding statement for a finitely presented module with a map. Supply the subsheaf K of the kernel used for injectivity and the finitely presented module G with a map used for surjectivity.
\end{itemize}\medskip
\ifdefined\hypertarget\hypertarget{sourcepointer-1}{}\fi
\noindent\textit{Source reference.} Original definition 002V: sources/stacks-chow/tags.txt, source lines 112--112. Prerequisite pointer invoked at proof line(s) 3267; it names a standard fact of a sibling chapter that the counted argument uses, and every step of the counted proof built on it is retained verbatim in the delivered body.\par\medskip
\ifdefined\hypertarget\hypertarget{sourcepointer-2}{}\fi
\noindent\textit{Source reference.} Original section 002D: sources/stacks-chow/tags.txt, source lines 94--94. Prerequisite pointer invoked at proof line(s) 3269; it names a standard fact of a sibling chapter that the counted argument uses, and every step of the counted proof built on it is retained verbatim in the delivered body.\par\medskip
\ifdefined\hypertarget\hypertarget{sourcepointer-3}{}\fi
\noindent\textit{Source reference.} Original reference 01PF: sources/stacks-properties/pinned-original-source.tex, source lines 3064--3064. Prerequisite pointer invoked at proof line(s) 3374; it names a standard fact of the same chapter that the counted argument uses, and every step of the counted proof built on it is retained verbatim in the delivered body.\par\medskip
\ifdefined\hypertarget\hypertarget{sourcepointer-4}{}\fi
\noindent\textit{Source reference.} Original reference 01PI: sources/stacks-properties/pinned-original-source.tex, source lines 3141--3141. Prerequisite pointer invoked at proof line(s) 3396; it names a standard fact of the same chapter that the counted argument uses, and every step of the counted proof built on it is retained verbatim in the delivered body.\par\medskip
\ifdefined\hypertarget\hypertarget{sourcepointer-5}{}\fi
\noindent\textit{Source reference.} Original lemma 01BC: sources/stacks-chow/tags.txt, source lines 1721--1721. Prerequisite pointer invoked at proof line(s) 3291; it names a standard fact of a sibling chapter that the counted argument uses, and every step of the counted proof built on it is retained verbatim in the delivered body.\par\medskip
\ifdefined\hypertarget\hypertarget{sourcepointer-6}{}\fi
\noindent\textit{Source reference.} Original section 009E: sources/stacks-chow/tags.txt, source lines 350--350. Prerequisite pointer invoked at proof line(s) 3368; it names a standard fact of a sibling chapter that the counted argument uses, and every step of the counted proof built on it is retained verbatim in the delivered body.\par\medskip
\begin{lemma}
\label{lemma-directed-colimit-diagram-finite-presentation}
\begin{slogan}
Quasi-coherent modules on quasi-compact and quasi-separated schemes
are filtered colimits of finitely presented modules.
\end{slogan}
Let $X$ be a scheme. Assume $X$ is quasi-compact and quasi-separated.
Let $\mathcal{F}$ be a quasi-coherent $\mathcal{O}_X$-module.
There exist
\begin{enumerate}
\item a filtered index category $\mathcal{I}$ (see
Categories, Definition \ref{categories-definition-directed}),
\item a diagram $\mathcal{I} \to \textit{Mod}(\mathcal{O}_X)$ (see
Categories, Section \ref{categories-section-limits}),
$i \mapsto \mathcal{F}_i$,
\item morphisms of $\mathcal{O}_X$-modules
$\varphi_i : \mathcal{F}_i \to \mathcal{F}$
\end{enumerate}
such that each $\mathcal{F}_i$ is of finite presentation
and such that the morphisms $\varphi_i$ induce an isomorphism
$$
\colim_i \mathcal{F}_i
=
\mathcal{F}.
$$
\end{lemma}

\begin{proof}
Choose a set $I$ and for each $i \in I$ an $\mathcal{O}_X$-module
of finite presentation and a homomorphism of $\mathcal{O}_X$-modules
$\varphi_i : \mathcal{F}_i \to \mathcal{F}$ with the following
property: For any $\psi : \mathcal{G} \to \mathcal{F}$ with $\mathcal{G}$
of finite presentation there is an $i \in I$ such that there exists
an isomorphism $\alpha : \mathcal{F}_i \to \mathcal{G}$ with
$\varphi_i = \psi \circ \alpha$. It is clear from
Modules, Lemma \ref{modules-lemma-set-isomorphism-classes-finite-type-modules}
that such a set exists (see also its proof).
We denote $\mathcal{I}$ the category
with $\Ob(\mathcal{I}) = I$ and given $i, i' \in I$
we set
$$
\Mor_\mathcal{I}(i, i') =
\{\alpha : \mathcal{F}_i \to \mathcal{F}_{i'} \mid
\alpha \circ \varphi_{i'} = \varphi_i
\}.
$$
We claim that $\mathcal{I}$ is a filtered category and that
$\mathcal{F} = \colim_i \mathcal{F}_i$.

\medskip\noindent
Let $i, i' \in I$. Then we can consider the morphism
$$
\mathcal{F}_i \oplus \mathcal{F}_{i'} \longrightarrow \mathcal{F}
$$
which is the direct sum of $\varphi_i$ and $\varphi_{i'}$.
Since a direct sum of finitely presented $\mathcal{O}_X$-modules
is finitely presented we see that there exists some $i'' \in I$
such that $\varphi_{i''} : \mathcal{F}_{i''} \to \mathcal{F}$
is isomorphic to the displayed arrow towards $\mathcal{F}$ above.
Since there are commutative diagrams
$$
\xymatrix{
\mathcal{F}_i \ar[r] \ar[d] & \mathcal{F} \ar@{=}[d] \\
\mathcal{F}_i \oplus \mathcal{F}_{i'} \ar[r] & \mathcal{F}
}
\quad
\text{and}
\quad
\xymatrix{
\mathcal{F}_{i'} \ar[r] \ar[d] & \mathcal{F} \ar@{=}[d] \\
\mathcal{F}_i \oplus \mathcal{F}_{i'} \ar[r] & \mathcal{F}
}
$$
we see that there are morphisms $i \to i''$ and $i' \to i''$
in $\mathcal{I}$. Next, suppose that we have $i, i' \in I$ and
morphisms $\alpha, \beta : i \to i'$ (corresponding to $\mathcal{O}_X$-module
maps $\alpha, \beta : \mathcal{F}_i \to \mathcal{F}_{i'}$).
In this case consider the coequalizer
$$
\mathcal{G} =
\Coker(
\mathcal{F}_i \xrightarrow{\alpha - \beta} \mathcal{F}_{i'}
)
$$
Note that $\mathcal{G}$ is an $\mathcal{O}_X$-module of finite presentation.
Since by definition of morphisms in the category $\mathcal{I}$
we have $\varphi_{i'} \circ \alpha = \varphi_{i'} \circ \beta$
we see that we get an induced map $\psi : \mathcal{G} \to \mathcal{F}$.
Hence again the pair $(\mathcal{G}, \psi)$ is isomorphic to
the pair $(\mathcal{F}_{i''}, \varphi_{i''})$ for some $i''$.
Hence we see that there exists a morphism $i' \to i''$ in
$\mathcal{I}$ which equalizes $\alpha$ and $\beta$. Thus we have
shown that the category $\mathcal{I}$ is filtered.

\medskip\noindent
We still have to show that the colimit of the diagram is $\mathcal{F}$.
By definition of the colimit, and by our definition of the category
$\mathcal{I}$ there is a canonical map
$$
\varphi :
\colim_i \mathcal{F}_i
\longrightarrow
\mathcal{F}.
$$
Pick $x \in X$. Let us show that $\varphi_x$ is an isomorphism.
Recall that
$$
(\colim_i \mathcal{F}_i)_x
=
\colim_i \mathcal{F}_{i, x},
$$
see
Sheaves, Section \ref{sheaves-section-limits-sheaves}.
First we show that the map $\varphi_x$ is injective.
Suppose that $s \in \mathcal{F}_{i, x}$ is an element
such that $s$ maps to zero in $\mathcal{F}_x$. Then there exists
a quasi-compact open $U$ such that $s$ comes from $s \in \mathcal{F}_i(U)$
and such that $\varphi_i(s) = 0$ in $\mathcal{F}(U)$.
By Lemma \ref{lemma-extend}
we can find a finite type quasi-coherent subsheaf
$\mathcal{K} \subset \Ker(\varphi_i)$ which restricts to
the quasi-coherent $\mathcal{O}_U$-submodule of $\mathcal{F}_i$
generated by $s$:
$\mathcal{K}|_U = \mathcal{O}_U\cdot s \subset \mathcal{F}_i|_U$.
Clearly, $\mathcal{F}_i/\mathcal{K}$ is of finite presentation and
the map $\varphi_i$ factors through the quotient map
$\mathcal{F}_i \to \mathcal{F}_i/\mathcal{K}$. Hence we can find
an $i' \in I$ and a morphism $\alpha : \mathcal{F}_i \to \mathcal{F}_{i'}$
in $\mathcal{I}$ which can be identified with the quotient map
$\mathcal{F}_i \to \mathcal{F}_i/\mathcal{K}$. Then it follows
that the section $s$ maps to zero in $\mathcal{F}_{i'}(U)$ and
in particular in
$(\colim_i \mathcal{F}_i)_x =
\colim_i \mathcal{F}_{i, x}$.
The injectivity follows.
Finally, we show that the map $\varphi_x$ is surjective.
Pick $s \in \mathcal{F}_x$. Choose a quasi-compact open neighbourhood
$U \subset X$ of $x$ such that $s$ corresponds to a section
$s \in \mathcal{F}(U)$. Consider the map
$s : \mathcal{O}_U \to \mathcal{F}$ (multiplication by $s$).
By Lemma \ref{lemma-extend-finite-presentation}
there exists an $\mathcal{O}_X$-module $\mathcal{G}$
of finite presentation and an $\mathcal{O}_X$-module map
$\mathcal{G} \to \mathcal{F}$ such that $\mathcal{G}|_U \to \mathcal{F}|_U$
is identified with
$s : \mathcal{O}_U \to \mathcal{F}$.
Again by definition of $\mathcal{I}$ there exists an $i \in I$
such that $\mathcal{G} \to \mathcal{F}$ is isomorphic to
$\varphi_i : \mathcal{F}_i \to \mathcal{F}$. Clearly there exists
a section $s' \in \mathcal{F}_i(U)$ mapping to $s \in \mathcal{F}(U)$.
This proves surjectivity and the proof of the lemma is complete.
\end{proof}

\end{document}
